\section{General activations and random-sign patterns}\label{sec:extensions}

The upper and lower bounds prove the Gaussian pressure formula for
$f\in C_c^\infty(\mathbb R)$. We now remove smoothness by estimates
uniform in the dimension and in the overlap path. For bounded activations,
Gaussian patterns permit approximation in Gaussian $L^1$ alone; the
random-sign extension also uses Gaussian-almost-everywhere continuity
to pass from discrete fields to Gaussian fields.
Let $\gamma$ denote standard Gaussian probability on
$\mathbb R$.

\subsection{Bounded Borel activations}

\begin{lemma}[Uniform approximation of the pattern functional]
\label{ext:path-approx}
If $f,g$ are bounded Borel functions with $|f|,|g|\le K$, then
\begin{align}
 \big|\mathbb E p_N(f)-\mathbb E p_N(g)\big|
 &\le\alpha e^{2K}\|f-g\|_{L^1(\gamma)},\label{ext:pressure-distance}\\
 \sup_{q\text{ step}}|V_f(q)-V_g(q)|
 &\le e^{2K}\|f-g\|_{L^1(\gamma)}.\label{ext:path-distance}
\end{align}
Consequently, for each bounded Borel $f$, the pattern functional $V_f$
extends uniquely and continuously in $L^1$ from step paths to $\mathcal Q$,
and the bound in \eqref{ext:path-distance} holds on all of $\mathcal Q$.
\end{lemma}

\begin{proof}
Interpolate linearly between $f$ and $g$. The derivative of the expected
pressure is a sum of Gibbs averages of $(g-f)(b_a(x))$, divided by $N$.
For each summand, remove pattern $a$. Restoring it changes a nonnegative
Gibbs average by a factor at most $e^{2K}$, because the removed log factor
lies in $[-K,K]$. Conditional on every other pattern and on a sampled spin
$x$, the fresh field $g^a\cdot x/\sqrt N$ has law $\gamma$. Hence the
absolute derivative is at most
$(M_N/N)e^{2K}\|g-f\|_{L^1(\gamma)}$. Integration proves
\eqref{ext:pressure-distance}.

For a fixed step path, terminal interpolation in the recursion propagates
$g-f$ by its tilted Gaussian kernels. This differentiation requires only
bounded terminal functions: it follows directly by differentiating the
finite expectations. In the notation of the pattern recursion, the density
of the entire tilted path relative to its independent Gaussian increments
has logarithm
\begin{equation}\label{ext:telescoped-density}
 \sum_{i=0}^{k}\zeta_i\big(U_i(Z_i)-U_{i-1}(Z_{i-1})\big)
 =f_t(Z_k)-\sum_{i=0}^{k-1}w_iU_i(Z_i),
 \qquad f_t=(1-t)f+tg.
\end{equation}
All the functions on the right take values in $[-K,K]$, so the density is
at most $e^{2K}$. Without the tilt, $Z_k$ is standard normal, since the
total variance is $q_k-q_{-1}=1$. This proves
\eqref{ext:path-distance}, uniformly over partitions and path values.

Choose $f_j\in C_c^\infty(\mathbb R)$ with $|f_j|\le K$ and
$\|f_j-f\|_{L^1(\gamma)}\to0$. Such approximations follow by spatial
truncation, approximation on compact intervals, and mollification with a
nonnegative smooth kernel. Each $V_{f_j}$ is continuous in $L^1$ by
Proposition~\ref{var:continuity-order}. Estimate~\eqref{ext:path-distance}
makes them uniformly Cauchy on the dense set of step paths and hence on
$\mathcal Q$. Their uniform limit is continuous and agrees with the
original finite recursion on every step path. Density also proves
uniqueness. The elementary bound $|V_f|\le\|f\|_\infty$ persists.
\end{proof}

\begin{proposition}[Completion of the bounded case]\label{ext:bounded}
Theorem~\ref{thm:main} holds for every bounded Borel activation.
\end{proposition}

\begin{proof}
Use the approximants from Lemma~\ref{ext:path-approx}. Since the entropy
functional is the same for every activation,
\[
 |\mathcal P_f(\alpha)-\mathcal P_{f_j}(\alpha)|
 \le\alpha\sup_{q\in\mathcal Q}|V_f(q)-V_{f_j}(q)|\longrightarrow0.
\]
Together with the uniform pressure estimate, the already proved smooth
formula gives convergence of $\mathbb E p_N(f)$ to
$\mathcal P_f(\alpha)$.

Changing a single whole pattern vector changes $\log Z_N(f)$ by at most
$2\|f\|_\infty$. The independent-input variance bound of
Lemma~\ref{arr:variance} therefore gives
\begin{equation}\label{ext:concentration}
 \operatorname{Var}(p_N(f))\le
 \frac{4M_N\|f\|_\infty^2}{N^2}=O(N^{-1}).
\end{equation}
Chebyshev's inequality proves convergence in probability.
\end{proof}

\subsection{An integrable unbounded class}

The factor $e^{2K}$ in the bounded approximation estimate is not useful
when the clipping level tends to infinity. Instead, we control the two
tails separately: integrability of $f$ controls the lower tail, while
integrability of $e^f$ controls the upper tail after a fixed lower clip.
This is why the limits below are taken in the stated order.

\begin{theorem}[Integrable activations]\label{ext:integrable}
Suppose $f:\mathbb R\to\mathbb R$ is Borel and
\begin{equation}\label{ext:moments}
 \mathbb E|f(G)|+\mathbb E e^{f(G)}<\infty.
\end{equation}
For $J,K>0$, let $f_{J,K}=\min\{K,\max\{-J,f\}\}$. The uniform limit
\begin{equation}\label{ext:unbounded-V}
 V_f=\lim_{J\to\infty}\lim_{K\to\infty}V_{f_{J,K}}
 \quad\text{on }\mathcal Q
\end{equation}
exists and is continuous in $L^1$. With this definition, the variational
formula \eqref{thm:formula} is finite and holds in expectation and in
probability.
\end{theorem}

\begin{proof}
We establish estimates uniform in both the dimension and the step path.
Let $P$ be the base probability for one fresh pattern, and let $t$ be
the value of $f$ on its field. In the spin model, $P$ is the Gibbs
measure with that whole pattern omitted, independent of the fresh row.
For a step path, let $(v_\lambda)$ be the cascade weights and let
$(y_\lambda)$ be the Gaussian field built from shared tree increments,
independent of those weights and with variance $q_{k-1}$ at each leaf.
Use the product probability and scalar function
\[
 P=\sum_\lambda v_\lambda\delta_\lambda\otimes\gamma,
 \qquad
 t(\lambda,\eta)=f\bigl(y_\lambda+\sqrt{1-q_{k-1}}\,\eta\bigr).
\]
Thus the residual $\eta$ is an integration coordinate of $P$, as in
Lemma~\ref{arr:fresh-pattern}. For a bounded Borel terminal function,
integrating $\eta$ gives the coefficient-one transform, and
Lemma~\ref{arr:tilt}, followed by the root Gaussian average, evaluates
the expected log integral as its pattern recursion.
This identity uses measurability and boundedness,
not the continuity needed for an array limit. Zero residual variance
and repeated covariance levels are allowed.
In the spin case, the fresh field at
each fixed state is standard normal. In the cascade case the same
marginal results after both averaging $y_\lambda$ and integrating
$\eta$. Independence of the weights and the shared field, followed by
Tonelli's theorem, therefore gives, for every nonnegative Borel $F$,
\begin{equation}\label{ext:fresh-marginal}
 \mathbb E P[F(t)]=\mathbb E F(f(G)).
\end{equation}
Here the expectation includes the independent pattern and any base data.
The log integral $\log P(e^t)$ is finite almost surely and integrable:
Jensen bounds it below by $P(t)$, while its positive part is bounded by
$P(e^t)$. These are integrable by \eqref{ext:moments} and
\eqref{ext:fresh-marginal}. Adding the rows successively also proves
integrability of the finite-volume pressure.

Put $t_J=\max(t,-J)$ and $r=(-J-t)_+$. The nonnegative lower-clipping
error satisfies
\begin{align}
 0\le \log P(e^{t_J})-\log P(e^t)
 &=-\log\frac{P(e^{t_J}e^{-r})}{P(e^{t_J})}
 \le\frac{P(e^{t_J}r)}{P(e^{t_J})}
 \le P(r).\label{ext:lower-clip}
\end{align}
The bound by the tilted mean of $r$ is Jensen's inequality under the
$e^{t_J}$ tilt. For the last inequality, $e^{\max(t,-J)}$ is increasing in
$t$, while $(-J-t)_+$ is decreasing. Their covariance under $P$ is
nonpositive. This follows also by expanding the expectation of
$(a(t)-a(t'))(b(t)-b(t'))\le0$ for two independent $P$ samples.

After the lower clip, the upper-clipped exponential is at least $e^{-J}$.
The inequality $\log(1+x)\le x$ gives
\begin{equation}\label{ext:upper-clip}
 0\le \log P(e^{t_J})-\log P(e^{\min(t_J,K)})
 \le e^J P[(e^t-e^K)_+].
\end{equation}
Define
\[
 \varepsilon_{J,K}=
 \mathbb E(-J-f(G))_+
 +e^J\mathbb E(e^{f(G)}-e^K)_+.
\]
Replacing one pattern at a time in the finite spin system, allowing the
other factors to be arbitrary existing clipped or unclipped factors, yields
\begin{equation}\label{ext:unbounded-pressure-error}
 \mathbb E|p_N(f)-p_N(f_{J,K})|\le\alpha\varepsilon_{J,K}.
\end{equation}
The omitted factor is independent of that base probability at each
replacement, so \eqref{ext:fresh-marginal} continues to apply. The same
argument for the single cascade log integral yields
\begin{equation}\label{ext:unbounded-path-error}
 \sup_{q\text{ step}}|V_f^{\mathrm{cas}}(q)-V_{f_{J,K}}(q)|
 \le\varepsilon_{J,K},
\end{equation}
where $V_f^{\mathrm{cas}}(q)$ denotes its integrable expected log value.

At fixed $J$, the upper-clipping bound compares each bounded recursion
uniformly on step paths with the cascade value for $\max(f,-J)$; its
error tends to zero as $K\to\infty$. Hence these bounded continuous
functionals are uniformly Cauchy on the dense set of step paths and
therefore on $\mathcal Q$. Their uniform limit is continuous. For
$J'\ge J$, apply the lower-clipping bound with $f$ replaced by
$\max(f,-J')$. The difference between the corresponding limits is at
most $\mathbb E(-J-f(G))_+$, which tends to zero as $J\to\infty$.
This proves the ordered uniform limit \eqref{ext:unbounded-V}. It agrees
with $V_f^{\mathrm{cas}}$ on step paths, and
\eqref{ext:unbounded-path-error}, with $V_f$ on the left, passes to all
paths by density.

Two applications of Jensen's inequality to the same cascade integral give
\[
 \mathbb E f(G)\le V_f(q)\le\log\mathbb E e^{f(G)}.
\]
The bounds hold first for step paths and then for every $q\in\mathcal Q$.
At $q=0$ the upper bound is an equality. Since $S_{\mathrm I}\ge0$ and
$S_{\mathrm I}(0)=0$, the variational value is finite, with
\[
 \alpha\mathbb E f(G)\le\mathcal P_f(\alpha)
 \le\alpha\log\mathbb E e^{f(G)}.
\]
Estimate~\eqref{ext:unbounded-path-error} also bounds the difference
between the clipped and unclipped variational infima by
$\alpha\varepsilon_{J,K}$.

Apply Proposition~\ref{ext:bounded} to $f_{J,K}$ and then
\eqref{ext:unbounded-pressure-error}. This proves convergence in
expectation. For convergence in probability, Markov's inequality controls
$p_N(f)-p_N(f_{J,K})$ uniformly in $N$; convergence for the bounded model
handles the middle term. Take $N\to\infty$, then $K\to\infty$ at fixed
$J$, and finally $J\to\infty$. No variance estimate for the unbounded
model is required.
\end{proof}

\begin{remark}[Why some activation restriction is needed]\label{ext:growth}
Wholly unrestricted real activations need not have finite limiting
pressure. For $f(s)=s^4$, choose $x_i=\operatorname{sign}(g_i^1)$. The law
of large numbers gives
$b_1(x)=N^{-1/2}\sum_i|g_i^1|\sim\sqrt{2/\pi}\sqrt N$ in probability.
All other pattern contributions are nonnegative, and the probability
weight of this one state is $2^{-N}$. Thus
$p_N(f)\ge b_1(x)^4/N-\log2\to\infty$ in probability.
\end{remark}

\subsection{Universality for random-sign patterns}

Disorder universality has also been established for broad hard-constraint
Ising perceptron classes by Nakajima and Sun~\cite{nakajimasun2022}.
The following result concerns finite bounded log-potentials and records
the continuity assumption used in the sign-to-Gaussian passage.

\begin{theorem}[Random-sign patterns]\label{ext:sign}
Replace the Gaussian pattern entries by independent uniform signs in
$\{-1,1\}$. If $f$ is bounded Borel and its discontinuity set has Gaussian
measure zero, then its expected pressure and its pressure in probability
converge to the same value $\mathcal P_f(\alpha)$ as in
Theorem~\ref{thm:main}. In particular, this holds for every bounded
continuous activation.
\end{theorem}

\begin{proof}
First let $f\in C_c^\infty(\mathbb R)$. We use the coordinate-replacement
form of the Lindeberg principle; see
Chatterjee~\cite[Theorem~1.1]{chatterjee2006lindeberg}.
Regard normalized pressure as a
function of a single pattern entry. Three differentiations of the finite
spin sum give a uniform bound $C_fN^{-5/2}$ for its third derivative.
Indeed, each derivative of the pattern argument supplies
$x_i/\sqrt N$, and the normalized log partition supplies the factor
$1/N$; the remaining terms are bounded Gibbs moments and cumulants of
$f',f'',f'''$. Taylor expansion through degree two and matching of the
first two Gaussian and sign moments therefore bound each entry replacement
by $C_fN^{-5/2}$. There are $M_NN=O(N^2)$ entries, so replacing them one
at a time gives an expected-pressure error $O_f(N^{-1/2})$.

For a general $f$ as in the statement, take smooth approximants $f_j$ with
$|f_j|,|f|\le K$ and $\|f_j-f\|_{L^1(\gamma)}\to0$. The omitted-pattern
argument from Lemma~\ref{ext:path-approx} still applies. Conditional on
any hypercube state, however, the fresh sign field now has the common law
\[
 X_N=N^{-1/2}\sum_{i=1}^N\epsilon_i,
 \qquad \mathbb P(\epsilon_i=1)=\mathbb P(\epsilon_i=-1)=1/2.
\]
Consequently,
\begin{equation}\label{ext:sign-approx}
 \big|\mathbb E p_N^{\mathrm{sign}}(f)
       -\mathbb E p_N^{\mathrm{sign}}(f_j)\big|
 \le\alpha e^{2K}\mathbb E|f(X_N)-f_j(X_N)|.
\end{equation}
The same third-order replacement calculation on smooth compactly
supported tests, together with $\mathbb EX_N^2=1$, proves
$X_N\Rightarrow G$. The bounded function $|f-f_j|$ is continuous outside
a Gaussian-null set. The usual bounded almost-everywhere-continuous
test criterion for weak convergence therefore gives
\[
 \mathbb E|f(X_N)-f_j(X_N)|
 \longrightarrow\mathbb E|f(G)-f_j(G)|.
\]
Take $N\to\infty$ in \eqref{ext:sign-approx} for fixed $j$, use the smooth
replacement result and the Gaussian formula, and then take $j\to\infty$.
This proves expected-pressure convergence. The bounded-pattern variance
bound \eqref{ext:concentration} holds for independent sign rows as well,
which proves convergence in probability.
\end{proof}

\begin{remark}[The regularity condition in sign universality]\label{ext:sign-scope}
One cannot replace the hypothesis of Theorem~\ref{ext:sign} by arbitrary
bounded Borel measurability. Take $f=\mathbf1_{\mathbb Q}$. For Gaussian
patterns, every field at every state is irrational almost surely, since
there are only finitely many states and rows at each size; hence $p_N=0$.
For sign patterns and square $N$, every field is rational, so
$p_N=M_N/N$. For odd nonsquare $N$, the numerator of each field is a
nonzero odd integer while $\sqrt N$ is irrational, so $p_N=0$.
The two subsequences give distinct limits. This function is discontinuous
everywhere and is excluded by Theorem~\ref{ext:sign}.
\end{remark}

\subsection{Normalization checks}

For $f=c$ constant, $V_f=c$ and $\inf_qS_{\mathrm I}(q)=0$, so
\eqref{thm:formula} gives $\alpha c$, as the partition function does
directly. A second check exercises both the diagonal shift and the duality.
For $f(s)=ts$, condition~\eqref{ext:moments} holds, and the finite Gaussian
recursion gives
\[
 V_f(q)=\frac{t^2}{2}\left(1-\int_0^1q(u)\,du\right).
\]
Set $H=\alpha t^2$. The constant field $h=H$ is an admissible trial, so
every variational objective is at least $H/2+\ell(H)$.
For the constant path $q_*=\mathbb E\tanh^2(\sqrt H\,G)$, the supporting
inequality \eqref{var:support} shows that the entropy supremum is attained
at that constant field. Equality follows, giving
\[
 \mathcal P_{ts}(\alpha)
 =H/2+\ell(H)=\mathbb E\log\cosh(\sqrt H\,G).
\]
This is also the direct answer: the linear activation produces independent
Gaussian coordinate fields with variances $t^2M_N/N\to H$, and the Ising
partition function factors into their $\cosh$ terms.
